\uuid{ASic}
\exo7id{4164}
\titre{exo7 4164}
\auteur{quercia}
\organisation{exo7}
\datecreate{2010-03-11}
\isIndication{false}
\isCorrection{true}
\chapitre{Fonction de plusieurs variables}
\sousChapitre{Dérivée partielle}
\cpge{mpsi-mp2i.fonctions-deux-variables.b, pcsi.fonctions-deux-variables.b, ptsi.fonctions-deux-variables.b, bcpst1.fonctions-deux-variables.c, bcpst1.developpements-limites.b, mp-mpi.calcul-differentiel-optimisation.b, pc.fonctions-plusieurs-variables.c, psi.fonctions-plusieurs-variables.c}
\module{Analyse}
\niveau{L2}
\difficulte{}

\contenu{
\question{
On considère la courbe d'équation $e^{x-y} = 1+2x+y$. Donner la tangente à cette
courbe et la position par rapport à la tangente au point $(0,0)$.
}
\reponse{
$y = -\frac x2 + \frac {9x^2}{16} + o (x^2)  \Rightarrow {}$
         tangente de pente $-\frac 12$, au dessus.
}
}
